\documentclass[12pt]{article}

\usepackage[utf8]{inputenc}
\usepackage{latexsym,amsfonts,amssymb,amsthm,amsmath}
\usepackage{tikz}
\usepackage{stanli}
\usetikzlibrary{shapes,calc,through,3d,patterns}
\usetikzlibrary{arrows.meta,decorations.pathreplacing,angles,quotes}
\usepackage{multicol}
\usepackage{geometry}
\usepackage{mathtools}
\usepackage{siunitx}
\usepackage{booktabs}

\setlength{\parindent}{0in}
\setlength{\oddsidemargin}{-0.25in}
\setlength{\textwidth}{7in}
\setlength{\textheight}{9.3in}
\setlength{\topmargin}{-1in}
\setlength{\headheight}{18pt}

\newcommand\blfootnote[1]{%
    \begingroup
    \renewcommand\thefootnote{}\footnote{#1}%
    \addtocounter{footnote}{-1}%
    \endgroup
}

\definecolor{darkred}{rgb}{0.6,0,0}

\newcommand{\bolddelta}{\boldsymbol\delta}
\newcommand{\boldxi}{\boldsymbol\xi}



\begin{document}

\noindent ME473 - Dynamic finite element analysis of structures\hfill EPFL\\
Stefano Burzio \hfill 2025\\
\vspace{0.3cm}

\hspace*{\fill}\textbf{\Large Problem set 7 - solutions}\hspace*{\fill}

\hrulefill

\subsection*{Problem 1}
\paragraph{a) Kronecker delta property:} Consider the shape functions associated with the AMC element, evaluated at the nodal coordinate $\boldsymbol{\xi}^j$. These are defined as
\[
    {}^a\mathbf{h}_i(\boldsymbol{\xi}^j) = \begin{bmatrix}
        (1 + \xi_1^i \xi_1^j)(1 + \xi_2^i \xi_2^j)\left(2 + \xi_1^i \xi_1^j + \xi_2^i \xi_2^j - (\xi_1^j)^2 - (\xi_2^j)^2\right)/8 \\
        b(1 + \xi_1^i \xi_1^j)(\xi_2^i + \xi_2^j)((\xi_2^j)^2 - 1)/8                                                               \\
        -a(\xi_1^i + \xi_1^j)((\xi_1^j)^2 - 1)(1 + \xi_2^i \xi_2^j)/8
    \end{bmatrix}^T.
\]
We shall verify the following:
\begin{itemize}
    \item if \(i = j\), the first component equals to 1, and the second and third components evaluate to 0;
    \item if \(i \neq j\), the first component equals to 0, and the second and third components evaluate to 0.
\end{itemize}
The nodal coordinates are given by:
\begin{center}
    \begin{tabular}{c|cc}
        $i$ & $\xi_1^i$ & $\xi_2^i$ \\\hline
        1   & $-1$      & $-1$      \\
        2   & $+1$      & $-1$      \\
        3   & $+1$      & $+1$      \\
        4   & $-1$      & $+1$
    \end{tabular}
\end{center}
To evaluate the first component of the shape function, we observe that $(1 + \xi_1^i \xi_1^j)(1 + \xi_2^i \xi_2^j)=4\delta_{ij}$. This can be verified through the following multiplication tables:
\begin{center}
    \begin{tabular}{c|cccc}
        $1+\xi_1^i\xi_1^j$ & $1$ & $2$ & $3$ & $4$ \\\hline
        1                  & $2$ & $0$ & $0$ & $2$ \\
        2                  & $0$ & $2$ & $2$ & $2$ \\
        3                  & $0$ & $2$ & $2$ & $2$ \\
        4                  & $2$ & $0$ & $0$ & $2$ \\
    \end{tabular}\qquad
    \begin{tabular}{c|cccc}
        $1+\xi_2^i\xi_2^j$ & $1$ & $2$ & $3$ & $4$ \\\hline
        1                  & $2$ & $2$ & $0$ & $0$ \\
        2                  & $2$ & $2$ & $0$ & $0$ \\
        3                  & $0$ & $0$ & $2$ & $2$ \\
        4                  & $0$ & $0$ & $2$ & $2$ \\
    \end{tabular}
\end{center}
Moreover, when $i=j$ the third factor of the first component simplifies to
$$2 + (\xi_1^i)^2 + (\xi_2^i)^2 - (\xi_1^j)^2 - (\xi_2^j)^2=2.$$
For the second and third components, note that they contain the factors $(\xi_2^j)^2 - 1$ and $(\xi_1^j)^2 - 1$, respectively, both of which vanish for all nodal values, as each $\xi_1^j=\xi_2^j = \pm1$. Therefore, these components are identically zero for all combinations of $i$ and $j$. Thus:
\[
    {}^a\mathbf{h}_i(\boldsymbol{\xi}^j) = \begin{bmatrix} \delta_{ij} \\ 0 \\ 0 \end{bmatrix}^T
\]
Notice that one can further show that
\[
    \frac{\partial{}^a\mathbf{h}_i}{\partial \xi_1}(\boldsymbol{\xi}^j) = \begin{bmatrix} 0\\ 0 \\ \delta_{ij}  \end{bmatrix}^T\qquad\text{and}\qquad\frac{\partial{}^a\mathbf{h}_i}{\partial \xi_2}(\boldsymbol{\xi}^j) = \begin{bmatrix} 0 \\ \delta_{ij} \\ 0 \end{bmatrix}^T.
\]
The Kronecker delta properties ensure that:
$${}^eu_3^h (\boldxi^j,t)=d^j,\qquad \partial_{\xi_1}{}^eu_3^h (\boldxi^j,t)=\theta_2^j,\qquad \text{ and} \qquad \partial_{\xi_2}{}^eu_3^h (\boldxi^j,t)=\theta_1^j.$$

\paragraph{b)Nonconformity of the AMC element:} when evaluating the shape functions of the AMC element along the side \( \xi_1 = +1 \),  the following expressions are obtained:
\begin{equation}
    \label{shape}
    {}^a\mathbf{h}_1 ={}^a\mathbf{h}_4 =
    \begin{bmatrix}
        0 \\
        0 \\
        0
    \end{bmatrix}^{T}
    \quad
    {}^a\mathbf{h}_2 =
    \begin{bmatrix}
        \frac{1}{4}(1-\xi_2)(2-\xi_2-\xi_2^2) \\
        (b/4)(-1+\xi_2)(\xi_2^2-1)            \\
        0
    \end{bmatrix}^{T}\quad   {}^a\mathbf{h}_3 =
    \begin{bmatrix}
        \frac{1}{4}(1+\xi_2)(2+\xi_2-\xi_2^2) \\
        (b/4)(1+\xi_2)(\xi_2^2-1)             \\
        0
    \end{bmatrix}^{T}.
\end{equation}
Consequently, the displacement field along this edge reads:
$$  {}^eu_3^h (1,\xi_2,t)  =\sum_{i=1}^4  {}^a\mathbf h_i(1,\xi_2)
    \begin{bmatrix}{}^ed^i(t) \\[2pt] {}^e{\theta}^i_1(t) \\[2pt] {}^e{\theta}^i_2(t)
    \end{bmatrix}.$$
It is clear from the form of the shape functions that the vertical displacement ${}^e u_3^h$ along the edge depends solely on the nodal values of the vertical displacements, ${}^e d^2$ and ${}^e d^3$, and rotations ${}^e \theta_1^2$ and ${}^e \theta_1^3$, associated with nodes 2 and 3, respectively. Therefore, the AMC element ensures the continuity of the vertical displacement along the side 2-3, satisfying the conditions of a $C^0$ element. (\emph{Caveat}: a similar verification should be conducted along side 3-4 to fully establish $C^0$ continuity.)

Furthermore, the rotations along the edge 2-3 ($\xi_1=1$) are given by:
\begin{align*}
    {}^e{\theta}_1^h(1,\xi_2,t) & =\frac{1}{b}\frac{\partial{}^eu_3^h}{\partial \xi_2}(1,\xi_2,t)=\frac{1}{b}\sum_{i=1}^4 \frac{\partial{}^a\mathbf{h}_i}{\partial \xi_2}(1,\xi_2) \begin{bmatrix}
                                                                                                                                                                                       {}^ed^i(t) \\[2pt] {}^e{\theta}^i_1(t) \\[2pt] {}^e{\theta}^i_2(t)
                                                                                                                                                                                   \end{bmatrix},   \\
    {}^e{\theta}_2^h(1,\xi_2,t) & =-\frac{1}{a}\frac{\partial{}^eu_3^h}{\partial \xi_1}(1,\xi_2,t)=-\frac{1}{a}\sum_{i=1}^4 \frac{\partial{}^a\mathbf{h}_i}{\partial \xi_1}(1,\xi_2) \begin{bmatrix}
                                                                                                                                                                                         {}^ed^i(t) \\[2pt] {}^e{\theta}^i_1(t) \\[2pt] {}^e{\theta}^i_2(t)
                                                                                                                                                                                     \end{bmatrix}.
\end{align*}
Considering the expressions of the shape functions \( {}^a\mathbf{h}_i \) along the side 2–3 provided in \eqref{shape}, it follows that \( {}^e \theta_1^h \) is uniquely determined by the nodal values \( {}^e d^2 \), \( {}^e d^3 \), \( {}^e \theta_1^2 \), and \( {}^e \theta_1^3 \). Therefore, continuity of the rotation \( {}^e \theta_1^h \) across the interface between adjacent elements is preserved (\emph{Caveat}: a similar verification should be conducted along side 3-4 to fully establish $C^0$ continuity.).

On the other hand, since
\[
    \frac{\partial {}^a\mathbf{h}_i}{\partial \xi_1}(1,\xi_2) = \begin{bmatrix}
        \frac{1}{8} (1+\xi_2^i\xi_2) \left( \xi_1^i(1+\xi_1^i+\xi_2^i\xi_2-\xi_2^2) + (1+\xi_1^i)(\xi_1^i-2) \right) \\[8pt]
        \frac{b}{8} \xi_1^i (\xi_2^i+\xi_2)(\xi_2^2-1)                                                               \\[8pt]
        \frac{-a}{4} (\xi_1^i+1)(1+\xi_2^i\xi_2)
    \end{bmatrix}^T,
\]
we have:
\begin{align*}
     & \frac{\partial {}^a\mathbf{h}_1^T}{\partial \xi_1} =
    \begin{bmatrix}
        \frac{1}{8}\xi_2(1-\xi_2^2)      \\
        -\frac{b}{8}(1+\xi_2)(\xi_2^2-1) \\
        0
    \end{bmatrix},
    \quad
     & \frac{\partial {}^a\mathbf{h}_2^T}{\partial \xi_1} =
    \begin{bmatrix}
        -\frac{1}{8}\xi_2(1-\xi_2^2)     \\
        \frac{b}{8}(-1+\xi_2)(\xi_2^2-1) \\
        -\frac{a}{2}(1-\xi_2)
    \end{bmatrix},                        \\
     & \frac{\partial {}^a\mathbf{h}_3^T}{\partial \xi_1} =
    \begin{bmatrix}
        \frac{1}{8}\xi_2(1-\xi_2^2)     \\
        \frac{b}{8}(1+\xi_2)(\xi_2^2-1) \\
        -\frac{a}{2}(1+\xi_2)
    \end{bmatrix},
     & \frac{\partial {}^a\mathbf{h}_4^T}{\partial \xi_1} =
    \begin{bmatrix}
        -\frac{1}{8}\xi_2(1-\xi_2^2)     \\
        -\frac{b}{8}(1+\xi_2)(\xi_2^2-1) \\
        0
    \end{bmatrix}.
\end{align*}
For ${}^e \theta_2^h$ to be continuous across the element interface, it would need to be uniquely determined by the degrees of freedom at nodes 2 and 3. However, the above expressions reveal that ${}^e \theta_2^h$ depends not only on ${}^e d^2$, ${}^e d^3$, ${}^e \theta_1^2$, ${}^e \theta_1^3$, ${}^e \theta_2^2$, and ${}^e \theta_2^3$, but also on the values of ${}^e d^j$ and ${}^e \theta_2^j$ at nodes $j=1$ and $j=4$. Consequently, the AMC element is classified as a \textbf{non-conforming element}: continuity of ${}^e \theta_2^h$ along the interface 2-3 is not guaranteed, and therefore the AMC element does not satisfy $C^1$ continuity requirements. By analogy, one can show that the continuity of ${}^e{\theta}_1^h$ fails along the interface given by the side 3-4. In spite of this, the AMC element is used and will, therefore, be considered further and the effect of this lack of continuity indicated.

\subsection*{Problem 2}
Consider the following ansatz for the transverse displacement approximation in a rectangular AMC element, characterized by four nodes and thus twelve degrees of freedom:
\begin{align*}
    {}^e u_3^h(x_1,x_2,t) = & \, a_1 + a_2 x_1 + a_3 x_2 + a_4 x_1^2 + a_5 x_1 x_2 + a_6 x_2^2                         \\
                            & + a_7 x_1^3 + a_8 x_1^2 x_2 + a_9 x_1 x_2^2 + a_{10} x_2^3 + a_{11} x_1^4 + a_{12} x_2^4
\end{align*}
Let us focus on the behavior along an edge of the element, for instance the bottom horizontal edge, defined by \( x_2 = 0 \). Along this edge, the displacement field simplifies to:
\[
    {}^e u_3^h(x_1,0,t) = a_1 + a_2 x_1 + a_4 x_1^2 + a_7 x_1^3 + a_{11} x_1^4
\]

When two adjacent elements share this edge, continuity of the displacement \( u_3^h \) across the common boundary requires that the polynomial expressions from both elements match along the entire edge. However, the presence of the quartic term \( x_1^4 \) implies that, beyond matching displacement and slope, the curvature \( \partial^2_{x_1 x_1} u_3^h \) would also need to be matched continuously.

Given that the AMC element formulation only introduces degrees of freedom associated with the displacement \( u_3^h \) and the rotations \( \theta_1^h \) and \( \theta_2^h \) at the nodes, it does not possess the necessary additional degrees of freedom to independently control second derivatives such as \( \partial^2_{x_1 x_1} u_3^h \). Consequently, continuity of curvature across the element boundaries cannot be enforced, leading to displacement discontinuities along shared edges.

In contrast, if the choice of the displacement function includes cross terms such as \( x_1^3 x_2 \) and \( x_1 x_2^3 \) instead of pure quartic terms, a different behavior is observed. Along the edge where \( x_2 = 0 \):  ${}^e u_3^h(x_1,0,t)$ is a cubic polynomial in $x_1$ ensuring that the displacement field \( u_3^h \) remains continuous across the edge without requiring higher-order derivative matching.

The inclusion of pure quartic terms \( x_1^4 \) and \( x_2^4 \) in the displacement approximation of the AMC plate bending element introduces uncontrolled curvature that cannot be matched at element interfaces, thereby resulting in displacement discontinuities. Conversely, the use of mixed cubic terms such as \( x_1^3 x_2 \) and \( x_1 x_2^3 \) enriches the internal displacement field while preserving continuity along edges, as their contributions naturally vanish along coordinate-aligned boundaries.

\subsection*{Problem 3}
Since ${}^e u_3^h$ is expressed as a linear combination of the shape functions ${}^a \mathbf{h}_i$, with time-dependent coefficients ${}^e d^i(t)$, ${}^e \theta^i_1(t)$ and ${}^e \theta^i_2(t)$ that are independent of the spatial variables, it follows that
\[
    \frac{\partial^2}{\partial \xi_1 \partial \xi_2} {}^e u_3^h(\boldsymbol{\xi},t) = \sum_{i=1}^4  \frac{\partial^2}{\partial \xi_1 \partial \xi_2}{}^a \mathbf{h}_i(\boldxi)
    \begin{bmatrix}
        {}^e d^i(t)        \\
        {}^e \theta_1^i(t) \\
        {}^e \theta_2^i(t) \\
    \end{bmatrix}.
\]
Accordingly, the problem reduces to computing the mixed second derivatives \( \frac{\partial^2}{\partial \xi_1 \partial \xi_2}{}^a \mathbf{h}_i\). Differentiating \({}^a \mathbf{h}_i(\boldsymbol{\xi})\) with respect to \(\xi_1\) yields
\[
    \frac{\partial }{\partial \xi_1}{}^a \mathbf{h}_i(\boldxi) = \begin{bmatrix}
        f_i^\prime(\xi_1)f_i(\xi_2)    \\
        bf_i^\prime(\xi_1)g_i(\xi_2)   \\
        -a g_i^\prime(\xi_1)f_i(\xi_2) \\
    \end{bmatrix}^T.
\]
Subsequently, differentiating with respect to $\xi_2$ leads to
$$
    \frac{\partial^2}{\partial \xi_1 \partial \xi_2}{}^a \mathbf{h}_i(\boldxi) = \begin{bmatrix}
        f_i^\prime(\xi_1)f_i^\prime(\xi_2)    \\
        bf_i^\prime(\xi_1)g_i^\prime(\xi_2)   \\
        -a g_i^\prime(\xi_1)f_i^\prime(\xi_2) \\
    \end{bmatrix}^T.
$$
At the nodal points, where \(\xi_1, \xi_2 \in \{\pm 1\}\), the derivatives of the Hermite functions are given by
\begin{align*}
    f_i^\prime(\xi) & = \frac{-3\xi^i (\xi^2 - 1)}{4},     \\
    g_i^\prime(\xi) & = \frac{3\xi^2 + 2\xi^i \xi - 1}{4}.
\end{align*}
Notably, $f_i^\prime(\pm 1) = 0$ since the factor $(\xi^2 - 1)$ vanishes at $\xi = \pm 1$. Thus, for every node $j=1,\dots,4$, the shape functions vanish:
\[
    \frac{\partial^2 }{\partial \xi_1 \partial \xi_2}{}^a \mathbf{h}_i(\boldxi^j)=\begin{bmatrix}
        0 \\
        0 \\
        0 \\
    \end{bmatrix}^T.
\]
Consequently,
\[
    \frac{\partial^2}{\partial \xi_1 \partial \xi_2} {}^e u_3^h(\boldsymbol{\xi}^j,t)= 0.
\]



\end{document}